\begin{musica}{Foi Deus}{Alberto Janes}{}
\begin{multicols}{2}
\ac{G}{Não} sei, \ac{C}{não} sabe nin\ac{G}{guém}
Por que canto o fado
Neste tom ma\ac{E}{goa}do
De dor e de \ac{Am}{pran}to.

E neste tor\ac{D7}{men}to
Todo o sofri\ac{D}{men}to,
Eu sinto que a \ac{Am}{a}lma
Cá dentro se a\ac{D7}{cal}ma
Nos versos que \ac{G}{can}to. \ac{D7}{}__

\ac{G}{Foi} Deus \ac{C}{que} deu luz aos \ac{G}{o}lhos,
Perfumou as rosas,
Deu oiro ao \ac{G7}{Sol} e prata ao \ac{C}{luar}.
Foi \ac{Cm}{Deus} que me pôs no \ac{G}{pei}to \ac{E7}{}__
Um rosário de \ac{Am}{pe}nas
Que vou des\ac{D7}{fian}do e choro a can\ac{G}{tar}.

\ac{D7}{E} pôs as estrelas no \ac{G}{céu},
\ac{E7}{E} fez o espaço sem \ac{Am}{fim},
\ac{F#}{Deu} o luto às ando\ac{Bm}{ri}nhas,
\ac{D7}{Ai}, e deu-me esta voz a \ac{G}{mim}. \ac{D7}{}__

Se canto, não sei o que canto
Misto de ventura, saudade, ternura
E talvez amor.

Mas sei que cantando
Sinto o mesmo quando
Se tem um desgosto,
E o pranto no rosto
Nos deixa melhor.

Foi Deus que deu voz ao vento,
Luz ao firmamento,
E deu o azul às ondas do mar.
Foi Deus que me pôs no peito
Um rosário de penas
Que vou desfiando e choro a cantar.

Fez poeta o rouxinol,
Pôs no campo o alecrim,
Deu as flores à Primavera
Ai, e deu-me esta voz a mim.
\end{multicols}

\posicoes{C: x32010, Cm: x35543, D: xx0232, D7: xx0212, E: 022100, E7: 020100,
  F#: 244322, G: 320033, G7: 320002, Am: x02210, Bm: x24432}
\end{musica}
